% SPDX-License-Identifier: CC-BY-SA-4.0
% Author: Matthieu Perrin

\begingroup

\begin{frame}{Asynchronisme}

  \begin{block}{Définition -- système asynchrone}
    \begin{itemize}
    \item Un système \alert{asynchrone} n'est pas régi par une horloge commune
    \item Pas de borne sur le nombre de pas d'un processus entre deux pas successifs d'un autre processus
    \end{itemize}
  \end{block}

  \begin{block}{Origines de l'asynchronisme}

    \begin{center}
      \begin{tikzpicture}[x=9mm, y=7mm]\small
        
        \draw[structure.fg] (0,5.3) node{Ordinateur};
        
        \draw[structure.fg] (9,5.3) node{Humain};
        
        \draw[structure.fg, -latex] (0,5) -- (0,0);
        \draw[structure.fg, -latex] (9,5) -- (9,0);
        
        \foreach \x in {0,0.5,...,3.5}{
          \draw[structure.fg!20] (-0.1,\x+1)    -- (9.1,\x+1);
          \draw[structure.fg!20] (-0.1,\x+0.85) -- (9.1,\x+0.85);
          \draw[structure.fg!20] (-0.1,\x+0.76) -- (9.1,\x+0.76);
          \draw[structure.fg!20] (-0.1,\x+0.7)  -- (9.1,\x+0.7);
          \draw[structure.fg!20] (-0.1,\x+0.65) -- (9.1,\x+0.65);
          \draw[structure.fg!20] (-0.1,\x+0.61) -- (9.1,\x+0.61);
          \draw[structure.fg!20] (-0.1,\x+0.577)-- (9.1,\x+0.577);
          \draw[structure.fg!20] (-0.1,\x+0.55) -- (9.1,\x+0.55);
          \draw[structure.fg!20] (-0.1,\x+0.52) -- (9.1,\x+0.52);
        }
        \foreach \x in {0,0.5,...,3.5}{
          \draw[structure.fg] (-0.1,\x+1)    -- (0.1,\x+1);
          \draw[structure.fg] (-0.1,\x+0.85) -- (0.1,\x+0.85);
          \draw[structure.fg] (-0.1,\x+0.76) -- (0.1,\x+0.76);
          \draw[structure.fg] (-0.1,\x+0.7)  -- (0.1,\x+0.7);
          \draw[structure.fg] (-0.1,\x+0.65) -- (0.1,\x+0.65);
          \draw[structure.fg] (-0.1,\x+0.61) -- (0.1,\x+0.61);
          \draw[structure.fg] (-0.1,\x+0.577)-- (0.1,\x+0.577);
          \draw[structure.fg] (-0.1,\x+0.55) -- (0.1,\x+0.55);
          \draw[structure.fg] (-0.1,\x+0.52) -- (0.1,\x+0.52);
        }
        \draw[structure.fg] (-0.1,0.5) -- (0.1,0.5);
        
        \foreach \x in {0,0.5,...,3.5}{
          \draw[structure.fg] (8.9,\x+1)    -- (9.1,\x+1);
          \draw[structure.fg] (8.9,\x+0.85) -- (9.1,\x+0.85);
          \draw[structure.fg] (8.9,\x+0.76) -- (9.1,\x+0.76);
          \draw[structure.fg] (8.9,\x+0.7)  -- (9.1,\x+0.7);
          \draw[structure.fg] (8.9,\x+0.65) -- (9.1,\x+0.65);
          \draw[structure.fg] (8.9,\x+0.61) -- (9.1,\x+0.61);
          \draw[structure.fg] (8.9,\x+0.577)-- (9.1,\x+0.577);
          \draw[structure.fg] (8.9,\x+0.55) -- (9.1,\x+0.55);
          \draw[structure.fg] (8.9,\x+0.52) -- (9.1,\x+0.52);
        }
        \draw[structure.fg] (8.9,0.5) -- (9.1,0.5);
        
        \draw (2.25,4.5) node{instruction};
        \draw (2.25,3) node{défaut de cache};
        \draw (2.25,1.5) node{accès au disque};
        \draw (2.25,0.5) node{granularité de l'ordonnanceur};
        
        \draw (6.75,4.5) node{battement de c\oe ur};
        \draw (6.75,3.1) node{douche};
        \draw (6.75,2.7) node{cours de fac};
        \draw (6.75,1.5) node{vacances de Noël};
        \draw (6.75,0.5) node{première mission martienne ?};
        
        \draw (0,4.5) node[left]{$10^{-9}$ s};
        \draw (0,3)   node[left]{$10^{-6}$ s};
        \draw (0,1.5) node[left]{$10^{-3}$ s};
        \draw (0,0.5) node[left]{$10^{-1}$ s};
        
        \draw (9,4.5) node[right]{$1$ s};
        \draw (9,3.625)   node[right]{$1$ min};
        \draw (9,2.75)   node[right]{$1$ h};
        \draw (9,2.04) node[right]{$1$ j};
        \draw (9,1.3) node[right]{$1$ mois};
        \draw (9,0.75) node[right]{$1$ an};
      \end{tikzpicture}
    \end{center}
    
  \end{block}
  
\end{frame}

\endgroup
\endinput
